\documentclass[11pt]{article}

\newcommand{\MemoTitle}{Don't Put Frontier AI Under ITAR: A Blunt Instrument for a Dual-Use Reality}
\newcommand{\MemoDate}{February 17, 2026}
\newcommand{\MemoLogoPath}{../../../tex/assets/icon-180.png}

% =========================================================
% Shared LaTeX preamble for all memos/*/tex/memo.tex
% =========================================================

\usepackage[letterpaper,top=1in,bottom=1in,left=1in,right=1in]{geometry}
\usepackage{microtype}
\usepackage[T1]{fontenc}
\usepackage[utf8]{inputenc}
\usepackage{lmodern}
\usepackage{graphicx}
\usepackage{xcolor}
\usepackage{hyperref}
\usepackage{xurl}
\usepackage{enumitem}

\setlength{\parindent}{0pt}
\setlength{\parskip}{0.35em}
\setlist[itemize]{leftmargin=1.2em,itemsep=0.04em,topsep=0.08em}
\setlist[enumerate]{leftmargin=1.4em,itemsep=0.04em,topsep=0.08em}

\definecolor{accentgold}{HTML}{D4A574}
\definecolor{ink}{HTML}{0F172A}

\hypersetup{
  hidelinks,
  pdftitle={\MemoTitle},
  pdfauthor={Adam Boas},
  pdfsubject={Policy Memo}
}

\pagestyle{empty}
\sloppy

\providecommand{\MemoTitle}{Memo}
\providecommand{\MemoDate}{\today}
\providecommand{\MemoLogoPath}{../../../tex/assets/icon-180.png}
\providecommand{\MemoAuthorName}{Adam Boas}
\providecommand{\MemoAuthorSite}{https://adamboas.com/}
\providecommand{\MemoAuthorSiteLabel}{adamboas.com}
\providecommand{\MemoAuthorEmail}{anboas@gmail.com}

\providecommand{\tightlist}{}

\newcommand{\MemoRenderHeader}{%
  \vspace*{-3.15em}
  \noindent
  \begin{minipage}[c]{0.14\textwidth}
  \raggedright
  \includegraphics[width=0.66\linewidth]{\MemoLogoPath}
  \end{minipage}%
  \begin{minipage}[c]{0.86\textwidth}
  \centering
  {\Large\bfseries\textcolor{ink}{\MemoAuthorName}}\\[-0.1em]
  {\small\textcolor{ink}{\href{\MemoAuthorSite}{\MemoAuthorSiteLabel} \quad|\quad \href{mailto:\MemoAuthorEmail}{\MemoAuthorEmail}}}\\[-0.1em]
  {\small\textcolor{ink}{\MemoDate}}
  \end{minipage}

  \vspace{0.35em}
  {\centering\large\bfseries \MemoTitle\par}
  \vspace{0.28em}
  \textcolor{accentgold}{\rule{\textwidth}{1.2pt}}
}


\begin{document}
\MemoRenderHeader

There is a seductive simplicity to the idea: if frontier AI models and
agentic runtimes can enable lethal autonomy, cyber operations, and
strategic-scale influence, then treat them like a weapon. Pull major AI
companies under the International Traffic in Arms Regulations (ITAR),
lock it down, and stop pretending it is ``just software.''

I am not convinced that works, and I am increasingly convinced it would
backfire.

ITAR is built to control defense articles, defense services, and
associated technical data on the U.S. Munitions List, administered by
the State Department. It is intentionally broad and punitive, and
``export'' includes disclosing controlled technical data to foreign
persons, including within the United States (a deemed export) {[}1{]}.
That is a feature for missiles and fire-control systems. It is a bug for
a technology ecosystem where capability is distributed across models,
data, tooling, evaluation harnesses, chips, and the global talent base
that builds and defends them.

\textbf{The innovation cost would not be a side effect; it would be the
main effect.}

If you ``ITARize'' frontier AI in a way that captures general-purpose
model development, you do not just throttle exports. You throttle basic
operating mechanics:

\begin{itemize}
\tightlist
\item
  \textbf{Hiring and collaboration become export-control exercises.} In
  ITAR-land, foreign-person access is not a rounding error; it is a
  compliance event. Day-to-day AI work, including code review, model
  debugging, red-teaming, and vendor incident response, depends on
  cross-border and cross-national participation. Treating routine
  collaboration as a regulated export will push companies into
  U.S.-person enclaves and split teams along citizenship lines {[}2{]}.
\item
  \textbf{Research disclosure becomes legally radioactive.} AI is not a
  single item you can box up. Competitiveness depends on iterative
  publication, open benchmarks, shared safety practices, and rapid
  diffusion of defensive techniques. Over-controlling ``technical data''
  risks converting a safety culture into a legal risk surface.
\item
  \textbf{Compliance becomes a moat.} ITAR overhead does not fall
  evenly. It advantages incumbents with deep compliance departments and
  punishes startups and research labs where genuine innovation often
  occurs. You can accidentally build a policy that ``secures'' AI by
  concentrating it.
\end{itemize}

This is not hypothetical. Even under the Commerce Department's Export
Administration Regulations (EAR)-based approach, which is typically more
flexible than ITAR, the U.S. has already seen backlash when proposed AI
diffusion controls were viewed as innovation- and diplomacy-hostile. In
May 2025, the Bureau of Industry and Security (BIS) described the
Biden-era AI Diffusion Rule as innovation-stifling and diplomatically
harmful, and rescinded it pending replacement {[}3{]}.

If that friction triggered a course correction under EAR, it is hard to
argue full ITAR treatment would be anything other than a strategic
self-inflicted wound.

\textbf{ITAR does not stop capability abroad; it mostly stops the U.S.
from participating.}

The uncomfortable truth is that frontier AI is already global. Export
controls can slow diffusion of specific chokepoints, such as chips,
select manufacturing tools, and some closed-weight artifacts, but cannot
prevent capable states and well-resourced firms from developing
comparable systems.

A maximalist ITAR approach would encourage:

\begin{itemize}
\tightlist
\item
  sovereign model programs outside U.S. jurisdiction,
\item
  offshore safety and evaluation ecosystems that evolve without U.S.
  participation,
\item
  and a bifurcated world where U.S. AI becomes compliance-heavy and
  slower-moving while others innovate around it.
\end{itemize}

That is not security. That is strategic isolation.

The U.S. government is already using a more targeted model. BIS has
explored controlling model weights for the most advanced closed-weight
AI models under EAR constructs. Congressional Research Service (CRS)
coverage of the AI Diffusion Rule proposal discussed restrictions on
model weights for advanced U.S. closed-weight dual-use AI models
{[}4{]}.

This is what scalable export control should look like: define a narrow
set of controllable artifacts and chokepoints while avoiding a sweep of
the domestic innovation base.

\textbf{The real problem is not ``AI exists''; it is ``AI acts without
governance''.}

The strongest ITARization argument is moral and operational: if AI can
be used for lethal autonomy at scale, we have crossed a line, so treat
enablers as munitions.

Even if you accept the premise, the conclusion does not follow.

The answer to ``Pandora's box is open'' is not ``regulate the box.'' It
is ``regulate the interfaces where the box touches reality.''

In practice, risk comes from systems that can:

\begin{itemize}
\tightlist
\item
  access sensitive data without enforced labeling and minimization,
\item
  invoke tools with side effects without authorization gates,
\item
  operate with unclear ``on behalf of'' semantics,
\item
  and produce actions without provenance.
\end{itemize}

Those are governance failures. They are solvable with architecture,
enforcement, and evidence. ITAR is an export regime, not an operational
control plane.

That is where an agent control-plane approach matters: enforceable
delegation, policy-gated action, and replayable operational evidence in
high-consequence environments (the Agent Control Plane Reference
Architecture, ACP-RA, is one proposed reference architecture for this).

\textbf{A smarter alternative.}

If the goal is reducing catastrophic misuse while preserving U.S.
innovation, a more defensible approach is:

\begin{enumerate}
\def\labelenumi{\arabic{enumi}.}
\tightlist
\item
  \textbf{Keep general-purpose frontier model development under EAR, not
  ITAR,} with narrowly defined controls on specific high-risk artifacts
  (for example: certain closed-weight model weights, specialized
  training/inference pipelines, and chip-enabled training services),
  updated as thresholds shift {[}4{]}.
\item
  \textbf{Apply ITAR only to defense-unique implementations,} including
  fine-tunes, toolchains, integrations, and services specifically
  designed for military end use (targeting, weapons employment, mission
  systems integration). This aligns with ITAR's purpose: regulating
  defense articles/services and direct technical data {[}1{]}.
\item
  \textbf{Require ``no identity, no action'' provenance for
  high-consequence agentic systems} in regulated environments (defense,
  critical infrastructure, classified or controlled domains). Make
  provenance operationally mandatory: signed delegation chains, policy
  decisions at gateways, and replayable evidence.
\end{enumerate}

This preserves what matters most: U.S. ability to innovate broadly,
collaborate with allies, and lead on safety, while still drawing hard
lines around defense-unique capabilities and sensitive artifacts. It
also gives leaders a practical bridge between policy and execution:
targeted export controls on one side, and operational governance on the
other.

\textbf{Conclusion.}

Putting ``AI companies'' under ITAR feels strong because ITAR is strong.
In this domain, strength is not measured by restrictiveness. It is
measured by whether outcomes beat what adversaries can get for free.

A sweeping ITAR move would slow U.S. innovation, fragment talent,
complicate allied collaboration, and push development offshore, while
doing little to stop capable actors elsewhere.

If we want to keep the future from becoming science fiction, we should
not start by regulating yesterday's way. We should regulate the
condition that actually creates harm: \textbf{uncontrolled autonomy
without enforceable identity, authorization, and action provenance.}

Control the interfaces where autonomy touches reality - identity,
authorization, and provenance - without freezing innovation.

\textbf{References.}

\begin{enumerate}
\def\labelenumi{\arabic{enumi}.}
\tightlist
\item
  International Traffic in Arms Regulations (ITAR) \textbar{} Wex
  \textbar{} LII.
  https://www.law.cornell.edu/wex/international\_traffic\_in\_arms\_regulations\_\%28itar\%29
\item
  International Traffic in Arms Regulations (ITAR) Compliance - Interim
  Policy \textbar{} University of Arkansas Research Integrity and
  Compliance. https://rsic.uark.edu/exportcontrol/itar-compliance.php
\item
  Department of Commerce announces rescission of Biden-era AI Diffusion
  Rule \textbar{} Bureau of Industry and Security.
  https://www.bis.gov/press-release/department-commerce-announces-rescission-biden-era-artificial-intelligence-diffusion-rule-strengthens
\item
  U.S. Export Controls and China: Advanced Semiconductors \textbar{} CRS
  \textbar{} Congress.gov. https://www.congress.gov/crs-product/R48642
\end{enumerate}

\end{document}
